(1) 定义 \(\omega_c=eB_0/m>0,\omega_p^2=ne^2/(m\epsilon_0)\)。从
\(-i\omega m\boldsymbol v=-e(\boldsymbol E+\boldsymbol v\times B_0\boldsymbol e_z)\)
得到 \(\boldsymbol v\times B_0\boldsymbol e_z=(B_0v_y,-B_0v_x,0)\)，故
\[
v_+=-{ie\over m(\omega+\omega_c)}E_+,\qquad
v_-=-{ie\over m(\omega-\omega_c)}E_-.
\]
(2) \(J_\pm=-ne v_\pm\)，用 \(J=-i\omega\epsilon_0(\epsilon_r-1)E\) 得
\[
\epsilon_\pm=1-{\omega_p^2\over\omega(\omega\pm\omega_c)},\qquad
c^2k_\pm^2=\omega^2\epsilon_\pm.
\]
截止 \(k=0\) 的正频根分别为
\[
\omega_{\mathrm{cut},+}={\sqrt{\omega_c^2+4\omega_p^2}-\omega_c\over2},\qquad
\omega_{\mathrm{cut},-}={\sqrt{\omega_c^2+4\omega_p^2}+\omega_c\over2}.
\]
负号支在 \(\omega=\omega_c\) 处有无碰撞回旋极点，应分开考察其两侧传播区；不能跨过极点连续使用无损响应。
(3) 由 \(c^2k^2=\omega^2-\omega\omega_p^2/(\omega\pm\omega_c)\) 对 \(\omega\) 求导得
\[
v_{g,\pm}={d\omega\over dk}=
{2c^2k\over 2\omega\mp\omega_p^2\omega_c/(\omega\pm\omega_c)^2}.
\]
该式仅对实传播 \(k\) 和非奇点有效。\(B_0\to0\) 两支合并 \(\epsilon=1-\omega_p^2/\omega^2,\ \omega^2=\omega_p^2+c^2k^2\)。
(4) 色散介质平均储能的电响应部分为 \(\epsilon_0[d(\omega\epsilon_\pm)/d\omega]|E_\pm|^2/4\)，且
\[
{d(\omega\epsilon_\pm)\over d\omega}
=1+{\omega_p^2\over(\omega\pm\omega_c)^2}>0 .
\]
再加正的磁场能，故总波能正定。局部 \(\epsilon_\pm<0\) 表示给定 \(\omega\) 无体传播波，不是物理负能量；模型失效于碰撞、离子共振及强非线性情况。
